% ===== PRÉAMBULE CANONIQUE — à coller VERBATIM en tête de CHAQUE .tex =====
\documentclass[11pt,a4paper]{article}
\usepackage[utf8]{inputenc}\usepackage[T1]{fontenc}\usepackage{lmodern}
\usepackage[french]{babel}
\usepackage{amsmath,amssymb,mathtools}\usepackage{icomma}
\usepackage{siunitx}\sisetup{locale=FR}  % \qty{}{}, \si{}, \ang{}
\usepackage{xcolor}\usepackage{colortbl}
\usepackage{tikz}\usepackage{pgfplots}\pgfplotsset{compat=1.18}
\usepackage[siunitx,europeanresistors]{circuitikz} % circuits (élec.)
\usepackage[most]{tcolorbox}\usepackage{enumitem}\usepackage{array,booktabs}
\usepackage[a4paper,margin=2.2cm]{geometry}
\usetikzlibrary{arrows.meta,calc,angles,quotes,positioning,patterns,%
  backgrounds,shapes.geometric,decorations.pathmorphing,decorations.markings,intersections}
\definecolor{primary}{RGB}{222,49,99}\definecolor{primarydark}{RGB}{150,22,55}
\definecolor{defcol}{RGB}{0,114,178}\definecolor{methcol}{RGB}{52,92,148}
\definecolor{excol}{RGB}{0,135,90}\definecolor{attncol}{RGB}{200,40,55}
\tcbset{boxbase/.style={enhanced,breakable,boxrule=0.9pt,arc=2.5pt,
  left=3mm,right=3mm,top=2.3mm,bottom=2.3mm,fonttitle=\bfseries,coltitle=white}}
% --- boîtes TUTORIEL (noms EXACTS, ne pas renommer) ---
\newtcolorbox{cle}[1][]{boxbase,colback=defcol!6,colframe=defcol,
  title={\textsc{À retenir}\ifx\relax#1\relax\relax\else\textmd{ : \itshape #1}\fi}}
\newtcolorbox{methode}[1][]{boxbase,colback=methcol!7,colframe=methcol,sharp corners,
  title={\textsc{Méthode}\ifx\relax#1\relax\relax\else\textmd{ --- \itshape #1}\fi}}
\newtcolorbox{exemplebox}[1][]{boxbase,colback=excol!5,colframe=excol,
  title={\textsc{Exemple}\ifx\relax#1\relax\relax\else\textmd{ : \itshape #1}\fi}}
\newtcolorbox{piege}[1][]{boxbase,colback=attncol!6,colframe=attncol,
  title={\textsc{Piège}\ifx\relax#1\relax\relax\else\textmd{ : \itshape #1}\fi}}
\newtcolorbox{remarque}[1][]{boxbase,colback=black!4,colframe=black!45,
  title={\textsc{Remarque}\ifx\relax#1\relax\relax\else\textmd{ : \itshape #1}\fi}}
% --- boîtes EXERCICE / CORRIGÉ (noms EXACTS, auto-comptées) ---
\newtcolorbox[auto counter]{exo}[1][]{boxbase,colback=primary!4,colframe=primary,
  title={\textsc{Exercice}~\thetcbcounter\ifx\relax#1\relax\relax\else\textmd{ --- \itshape #1}\fi}}
\newtcolorbox[auto counter]{corrige}[1][]{boxbase,colback=excol!4,colframe=excol,
  title={\textsc{Corrigé}~\thetcbcounter\ifx\relax#1\relax\relax\else\textmd{ --- \itshape #1}\fi}}
\newcommand{\vv}[1]{\vec{#1}}\newcommand{\ds}{\displaystyle}
\newcommand{\R}{\mathbb{R}}\newcommand{\N}{\mathbb{N}}\newcommand{\Z}{\mathbb{Z}}
% (un \usepackage{fancyhdr}+en-tête est possible ; il est ignoré côté web)
% =========================================================================
\begin{document}
\sisetup{per-mode=symbol}

\begin{exo}[théorème avec deux vitesses non nulles]
Un wagonnet de masse $m=\qty{200}{\kilogram}$ voit sa vitesse passer de $v_A=\qty{2}{\metre\per\second}$
à $v_B=\qty{5}{\metre\per\second}$. Calcule le travail total $\sum W$ reçu par le wagonnet entre $A$
et $B$.
\end{exo}

\begin{exo}[distance de freinage]
Une voiture de masse $m=\qty{1200}{\kilogram}$ roule à $v=\qty{20}{\metre\per\second}$. En freinant,
les roues exercent une force résistante constante $F=\qty{6000}{\newton}$ jusqu'à l'arrêt.
\begin{enumerate}[label=\alph*)]
  \item Écris le théorème de l'énergie cinétique entre le début du freinage et l'arrêt.
  \item Déduis-en la distance de freinage $d$.
\end{enumerate}
\end{exo}

\begin{exo}[force de frottement à partir de la variation d'énergie cinétique]
Un palet de masse $m=\qty{0,5}{\kilogram}$, lancé à $v=\qty{4}{\metre\per\second}$ sur une table
horizontale, s'arrête après $d=\qty{10}{\metre}$ à cause des frottements.
\begin{enumerate}[label=\alph*)]
  \item Quelle est la variation d'énergie cinétique $\Delta E_c$ entre le lancement et l'arrêt ?
  \item Ce travail est celui de la force de frottement ($-F_f\,d$). Déduis-en l'intensité $F_f$.
\end{enumerate}
\end{exo}

\begin{exo}[l'énergie cinétique ne dépend pas de la pesanteur]
Un objet de masse $m=\qty{3}{\kilogram}$ se déplace horizontalement à $v=\qty{10}{\metre\per\second}$.
\begin{enumerate}[label=\alph*)]
  \item Calcule son énergie cinétique sur la Terre ($g=\qty{9,81}{\metre\per\second\squared}$).
  \item Sur la Lune, où la pesanteur vaut environ $g/6$, l'objet se déplace à la \emph{même} vitesse
        $v=\qty{10}{\metre\per\second}$. Quelle est son énergie cinétique ? Justifie.
\end{enumerate}
\end{exo}

\begin{exo}[la distance de freinage varie comme le carré de la vitesse]
À $\qty{50}{\kilo\metre\per\hour}$, une voiture s'arrête sur une distance $d_1=\qty{25}{\metre}$
(force de freinage constante).
\begin{enumerate}[label=\alph*)]
  \item En freinant avec la \emph{même} force, la distance d'arrêt est proportionnelle à $v^2$.
        Écris le rapport $\dfrac{d_2}{d_1}$ en fonction des vitesses.
  \item Déduis-en la distance d'arrêt $d_2$ si la voiture roule à $\qty{70}{\kilo\metre\per\hour}$.
\end{enumerate}
\end{exo}
\end{document}
